% from ega4-2-fr.tex
\oldpage[IV-2]{166}
contient un ouvert non vide $V'_1$; comme $g$ est fermé (II, 6.1.10), on
peut supposer que $V'_1=g^{-1}(V')$, où $V'$ est un ouvert affine de $X'$;
remplaçant $A'$ par l'anneau de $V'$ et $A'_1$ par celui de $V'_1$, on peut
donc supposer que $X'_1$ est régulier. En outre, le même raisonnement que
dans I) appliqué au point générique $\xi'$ de $X'$ permet (en remplaçant au
besoin $X'$ par un voisinage de $\xi'$) de supposer que $A'_1$ est un
$A'$-module libre, et par suite que le morphisme $g$ est plat. Mais il
résulte alors de (6.5.2, (i)) que $X'$ est régulier. C.Q.F.D.
\end{proof}

\begin{corollary}[6.12.5]
\label{IV.2.s6.12.5-fr}
\textnormal{(Zariski).} --- Soit $k$ un corps; pour tout $k$-préschéma $X$
localement de type fini sur $k$, l'ensemble des $x\in X$ où $X$ est
régulier (resp. géométriquement régulier sur $k$) est ouvert dans $X$.
\end{corollary}

\begin{proof}
L'assertion du corollaire relative à la propriété d'être régulier résulte de
(6.12.4) en prenant $A=k$. L'assertion relative à la propriété d'être
géométriquement régulier résulte déjà de (6.8.7); on peut aussi la déduire de
(6.12.4) de la façon suivante: posons $k'=k^{p^{-\infty}}$ ($p$ exposant
caractéristique de $k$); comme le morphisme
$\operatorname{Spec}(k')\to\operatorname{Spec}(k)$ est radiciel, entier et
surjectif, c'est un homéomorphisme universel (2.4.5), donc le morphisme
projection $X\otimes_kk'\to X$ est un homéomorphisme. La projection dans
$X$ de $\operatorname{Reg}(X\otimes_kk')$ est l'ensemble des points de $X$
où $X$ est géométriquement régulier sur $k$, en vertu de (6.7.7, e)), donc
cet ensemble est ouvert d'après ce qui précède.
\end{proof}

\begin{corollary}[6.12.6]
\label{IV.2.s6.12.6-fr}
Soit $A$ un anneau ayant l'une des propriétés suivantes:
\begin{enumerate}
\item[(i)] $A$ est un anneau de Dedekind et son corps des fractions $K$ est
de caractéristique 0.
\item[(ii)] $A$ est un anneau semi-local noethérien de dimension $\leq1$.
\end{enumerate}
Alors, pour tout préschéma $X$ localement de type fini sur $A$,
$\operatorname{Reg}(X)$ est ouvert dans $X$.
\end{corollary}

\begin{proof}
Vérifions dans les deux cas la condition c) de (6.12.4). Dans les deux cas,
un idéal premier $\mathfrak p$ de $A$ est maximal ou minimal; si
$\mathfrak p$ est maximal, une $(A/\mathfrak p)$-algèbre finie et intègre
est un corps et la condition c) de (6.12.4) est trivialement vérifiée.
Supposons donc $\mathfrak p$ non maximal, et distinguons les deux cas de
l'énoncé.

(i) Si $K$ est de caractéristique 0, il n'y a pas d'autre extension
radicielle de $K$ que $K$ lui-même; comme un anneau de Dedekind est régulier
(0, 17.1.4), la condition c) de (6.12.4) est trivialement vérifiée.

(ii) On peut alors supposer $A$ intègre (6.12.2), soit $K$ son corps des
fractions; si $K'$ est une extension radicielle finie de $K$, $A'$ une
sous-$A$-algèbre de $K'$ engendrée par un système fini de générateurs de
$K'$ sur $K$, entiers sur $A$, $A'$ est un anneau semi-local intègre de
dimension 1 (0, 16.1.5), et par suite, dans
$X'=\operatorname{Spec}(A')$, l'ensemble réduit au point générique est
ouvert et évidemment contenu dans $\operatorname{Reg}(X')$, ce qui prouve
la condition c) de (6.12.4) dans ce cas.
\end{proof}

Ce corollaire s'applique notamment lorsque $A=\mathbf Z$.

\begin{theorem}[6.12.7]
\label{IV.2.s6.12.7-fr}
\textnormal{(Nagata).} --- Soient $A$ un anneau local noethérien complet,
$X=\operatorname{Spec}(A)$. Alors $\operatorname{Reg}(X)$ est ouvert dans
$X$.
\end{theorem}

\begin{proof}
Compte tenu de (6.12.2), on est ramené au cas où de plus $A$ est intègre, et
à prouver que dans ce cas $\operatorname{Reg}(X)$ contient une partie
ouverte non vide de $X$. Distinguons deux cas:

